\documentclass{article}

\usepackage[margin = 1in]{geometry}

\pagestyle{empty}

\begin{document}

Name: \underline{\phantom{xxxxxxxxxxxxxxxxxxxx}}

Principles of Microeconomics

Reading Check: Ch. 4 Intro - 4.2.

\begin{enumerate}

\item A competitive market is a market in which

	\begin{enumerate}
	
	\item an auctioneer helps set prices and arrange sales.
	
	\item there are only a few sellers.
	
	\item the forces of supply and demand do not apply.
	
	\item no individual buyer or seller has any significant impact on the market price.
	
	\end{enumerate}
	
\end{enumerate}

\bigskip

Name: \underline{\phantom{xxxxxxxxxxxxxxxxxxxx}}

Principles of Microeconomics

Reading Check: Ch. 4 Intro - 4.2.

\begin{enumerate}

\item Which of the following changes would \textit{not} shift the demand curve for a good or service?

	\begin{enumerate}
	
	\item A change in income.
	
	\item A change in the price of the good or service.
	
	\item A change in expectations about the future price of the good or service.

	\item A change in the price of a related good or service.
	
	\end{enumerate}
	
\end{enumerate}

\bigskip

Name: \underline{\phantom{xxxxxxxxxxxxxxxxxxxx}}

Principles of Microeconomics

Reading Check: Ch. 4 Intro - 4.2.

\begin{enumerate}

\item If the demand for a good falls when income falls, then the good is called

	\begin{enumerate}
	
	\item a normal good.
	
	\item a regular good.
	
	\item an ordinary good.
	
	\item an inferior good.
	
	\end{enumerate}

\end{enumerate}

\bigskip

Name: \underline{\phantom{xxxxxxxxxxxxxxxxxxxx}}

Principles of Microeconomics

Reading Check: Ch. 4 Intro - 4.2.

\begin{enumerate}

\item Which of the following demonstrates the law of demand?

	\begin{enumerate}
	
	\item After Wyatt got a raise at work, he bought more pretzels at \$1.50 per bag than he did before his raise.
	
	\item Sanjay buys fewer pastries at \$0.75 per pastry than at \$1 per pastry, other things equal.
	
	\item Ellie buys more apple fritters at \$0.25 per apple fritter than at \$0.50 per apple fritter, other things equal. 
	
	\item Trinity buys fewer candy bars at \$0.60 after the price of Skittles falls to \$0.50 per bag.
	
	\end{enumerate}

\end{enumerate}

\newpage

Name: \underline{\phantom{xxxxxxxxxxxxxxxxxxxx}}

Principles of Microeconomics

Reading Check: Ch. 4.3. - 4.4.

\begin{enumerate}

\item Which of the following events would cause a movement upward and to the right along the supply curve for kiwis?

	\begin{enumerate}
	
	\item The number of sellers of kiwis increases.

	\item There is an advance in technology that reduces the cost of producing kiwis.

	\item The price of water decreases, and water is an input in the production of kiwis.
	
	\item The price of kiwis rises.
	
	\end{enumerate}

\end{enumerate}

\bigskip

Name: \underline{\phantom{xxxxxxxxxxxxxxxxxxxx}}

Principles of Microeconomics

Reading Check: Ch. 4.3. - 4.4.

\begin{enumerate}

\item If the supply of a product increases, then we would expect equilibrium price

	\begin{enumerate}
	
	\item to increase and equilibrium quantity to decrease.
	
	\item to decrease and equilibrium quantity to increase.
	
	\item and equilibrium quantity to both increase.
	
	\item and equilibrium quantity to both decrease.
	
	\end{enumerate}

\end{enumerate}

\bigskip

Name: \underline{\phantom{xxxxxxxxxxxxxxxxxxxx}}

Principles of Microeconomics

Reading Check: Ch. 4.3. - 4.4.

\begin{enumerate}

\item The market supply curve

	\begin{enumerate}
	
	\item shows how supply changes when expectations change.
	
	\item slopes downward.
	
	\item represents the sum of the prices that all the sellers are willing to accept for a given quantity of the good.
	
	\item represents the sum of the quantities supplied by all the sellers at each price of the good. 
	
	\end{enumerate}

\end{enumerate}

\bigskip

Name: \underline{\phantom{xxxxxxxxxxxxxxxxxxxx}}

Principles of Microeconomics

Reading Check: Ch. 4.3. - 4.4.

\begin{enumerate}

\item If the demand for a product increases, then we would expect equilibrium price

	\begin{enumerate}
	
	\item to increase and equilibrium quantity to decrease.
	
	\item to decrease and equilibrium quantity to increase.
	
	\item and equilibrium quantity both to increase.
	
	\item and equilibrium quantity both to decrease.
	
	\end{enumerate}

\end{enumerate}

\end{document}